% titel: Blatt 6, A4: Überlagerte Schwingungen
% skript: Kap. 1 (komplexe Darstellung, Zeiger)

% Teil I: Werkzeuge, kein Lösungsweg.

Zwei gleichfrequente Kosinus-Schwingungen zu einer Schwingung
$A\cos(\omega t+\varphi)$ zusammenfassen.

\begin{werkzeugkasten}[Phasor/Zeiger-Methode]
Einen Kosinus $a\cos(\omega t+\varphi)$ deutet man als Realteil des
Zeigers $a\mathrm e^{i\varphi}\cdot\mathrm e^{i\omega t}$. Die Überlagerung
gleicher Frequenz ist die Summe der Zeiger (eine komplexe Zahl):
\begin{equation*}
  \sum_j a_j\cos(\omega t+\varphi_j)
  =\operatorname{Re}\Bigl[\Bigl(\sum_j a_j\mathrm e^{i\varphi_j}\Bigr)
  \mathrm e^{i\omega t}\Bigr]
  =A\cos(\omega t+\varphi),
\end{equation*}
wobei $A$ der Betrag und $\varphi$ das Argument der Zeigersumme ist.

\textit{Skript:} Kap. 1 (komplexe Zahlen in Schwingungen).
\end{werkzeugkasten}

\begin{werkzeugkasten}[Schrittfolge]
\begin{enumerate}
  \item Jede Schwingung als Zeiger schreiben:
        $\cos(\omega t)\mapsto 1\cdot e^{i0}$, $2\cos(\omega t+2\pi/3)
        \mapsto2\,e^{i2\pi/3}$.
  \item Zeiger addieren und in Polarkoordinaten umrechnen
        (hier Ergebnis $i\sqrt3=\sqrt3\,e^{i\pi/2}$).
  \item Ablesen: $A=\sqrt3$, $\varphi=\pi/2$.
\end{enumerate}
\end{werkzeugkasten}

\begin{achtung}[Darauf achten]
\begin{itemize}
  \item Der Zeiger erfasst nur Amplitude und Phase; die gemeinsame
        $e^{i\omega t}$-Zeitabhängigkeit wird ausgeklammert.
  \item Das Ergebnis $i\sqrt3$ hat Realteil $0$: die Phasenverschiebung ist
        $\varphi=\pi/2$, die Amplitude $\sqrt3$, \emph{nicht} $0$.
  \item Nach dem Umrechnen die Form $A\cos(\omega t+\varphi)$ mit $A>0$
        und $\varphi\in[0,2\pi)$ aufschreiben.
  \item Verifikation an einer Zeitachse (z.B.\ $\omega t=0$) hilft, das
        Vorzeichen der Phase zu prüfen.
\end{itemize}
\end{achtung}
